\chapter{System Requirements and Overall Architecture}


\section{Requirements Analysis}

The system is designed for day-to-day programming-course teaching and primarily serves students and instructors. Students mainly care about smooth submission, clear feedback, and timely results, while instructors care more about convenient class/assignment management, stable grading behavior, and traceable workflows.

Based on practical usage scenarios, the system objectives can be summarized as follows:

\begin{enumerate}
\item Support a complete closed loop from submission to grading feedback.
\item Ensure grading results are interpretable and reviewable.
\item Support continuous instructor-side management of classes, students, and assignments.
\item Support online deployment and off-campus access for home-use scenarios.
\end{enumerate}


\paragraph{Student-side Functional Requirements}
Core student requirements include:

\begin{enumerate}
\item Account login and student registration.
\item Assignment list viewing and code submission.
\item Result viewing, including runtime output, static-analysis score, final score, and AI feedback.
\item Basic profile maintenance (e.g., name).
\end{enumerate}


\paragraph{Instructor-side Functional Requirements}
In addition to score review, instructors require complete management capabilities:

\begin{enumerate}
\item Instructor login and role-based authentication.
\item Class management: create, rename, delete, add/remove students.
\item Student management: view student details, review submission records, delete abnormal accounts.
\item Assignment-view dashboard: review student status and scores by assignment.
\item Report viewing: jump directly to detailed submission reports.
\item Grade export: CSV export for offline statistics.
\end{enumerate}


\paragraph{Non-functional Requirements}
Besides business functions, the system should satisfy:

\begin{enumerate}
\item Stability: maintain service availability under repeated submission scenarios.
\item Security: execute untrusted code in controlled environments rather than host runtime.
\item Consistency: keep scoring rules reproducible and avoid large fluctuations across similar submissions.
\item Scalability: support future expansion to more problems, languages, and finer analysis rules.
\item Deployability: support Docker-based deployment and public-access configuration.
\end{enumerate}


\section{Overall System Architecture}

The system adopts a front-end/back-end separated architecture with four layers:

\begin{enumerate}
\item Presentation layer: student and instructor interfaces for interaction and result display.
\item Service layer: authentication, submission processing, scoring feedback, and management APIs.
\item Data layer: PostgreSQL persistence for users, assignments, submissions, feedback, and classes.
\item Retrieval/feedback layer: AI feedback generation grounded in course-document retrieval.
\end{enumerate}


At the deployment layer, services are orchestrated by Docker Compose. The frontend uses Nginx to proxy backend APIs, enabling quick startup in local or public environments.

\section{Core Workflow Design}


\paragraph{Submission Workflow}
After code submission, backend processing follows these steps:

\begin{enumerate}
\item Validate assignment and user identity.
\item Store submission version and code content.
\item Execute code via execution service.
\item Generate issue list and static score via static-analysis module.
\item Compute final score via grading module.
\item Generate RAG-based feedback and write it to the database.
\end{enumerate}


This workflow ensures each submission is traceable and reviewable.

\paragraph{Scoring Workflow}
A two-track strategy is used:

\begin{enumerate}
\item Final score: based mainly on hidden test outcomes, reflecting functional correctness.
\item Static-analysis score: reflecting code quality and implementation process, without replacing the final score.
\end{enumerate}


To prevent inflated scores from template-only submissions, blank-submission detection is applied. If only template content or no effective new code is detected, static-analysis scores are significantly reduced or set to zero, and a warning is shown on the result page.

\paragraph{Instructor Management Workflow}
The instructor workbench links class, assignment, and student views. It supports class-based filtering, assignment-based review, and student-level drill-down to detailed submission reports, integrating score review and teaching management in one interface workflow.

